% r03_h1s_c: proof of the two-flow comparison bound, part 3: degree bookkeeping per ordering after the rule; the enumeration.

\begin{lemma}[degree bookkeeping]\label{rot:h1s:Dp}
Fix an all-$C$ component $G$ of a monomial, a derivative placement, an ordering of its scale-summed elements, and apply (NF1)--(NF4) and (R1)--(R3). Then in the reduced ordering every connected component $H$ of every initial segment satisfies
\[
D(H)\ge1,\quad\text{or}\quad D(H)=n+m\alpha\ \ (n\ge0,\ m\ge1;\ H\text{ contains a compound on an external leg}),\quad\text{or}\quad H\text{ is an explicit-$\eta'$ component}
\]
(a $Y$-displacement term of (R2b) or a $\partial^{\eta'}g$ term of (R2a), with $D(H)=1-\theta$ after the gain $\eta'\ge\eta'^\theta$ is taken out), with at most $n_0\le\bar n$ zero-degree components in the per-flow count (local vertices of weight $0$, explicit-$\eta'$ components). In particular $D(H)-\theta\ge\min(\alpha,1)-\theta>0$ for $\theta\le\theta_1$ whenever $H$ contains the weakened element.
\end{lemma}
\begin{proof}
(a) Degrees are a property of the multilinear form and flow-independent: both flows have the same lines, the same placements and the same transfers, since (NF2) and the rule are applied in each flow's own frame to the same combinatorial objects.

(b) \emph{No component has $D<0$.} By \Cref{rot:h1s:CL} a connected line set with $D_0<0$ is a unit $U$. In the reduced ordering $U$'s lines are either inside a compound, and then $U$ alone is not an element, or $U$ is split, and then both legs $f$, $g$ of $U$ are earlier than $\sigma_{\max}(U)$: the first initial segment containing both lines of $U$ contains $f$ and $g$, attached at $U$'s vertices, so the component is $\{f,g,U,\dots\}$ with $D\ge d_f+d_g-4+9\ge1$ ($d_f,d_g\ge-2$ by \Cref{rot:h1s:nd}(i)); the segment containing only $U$'s first line, the lone $\partial\partial C$ of a split self-energy with $D=0$ when it precedes both legs, is a degree-zero component covered by (c). Transfers never create $D<0$: a transfer lands on an element later than its source component $H_0$ and attached to a vertex of $H_0$, so every later component containing the recipient contains $H_0$, where the $+1$ and the $-1$ cancel (this is \cite[(2.8)--(2.9)]{Ba3}: ``the degrees of these graphs are $\ge D(G_0)$''); the components containing $H_0$ but not the recipient rose by $1$; no other component changed. Loop closure after a transfer is harmless for the same reason: if the recipient line later closes a loop, the component $H\cup l$ has $D(H\cup l)\ge D_0(H\cup l)\ge-1$, with $-1$ only for a bare unit; then $H$ is the first line of a split unit and $l$ its other line, the coarsest, so both legs are earlier and attached when $l$ enters, and the component is $\{f,g,U,\dots\}$ with $D\ge1$ as above.

(c) \emph{Components with $D(H)=0$ at formation.} Since every transfer inside $H$ with recipient inside $H$ is neutral and one with recipient outside is $+1$, $D(H)\ge D_0(H)$, so $D_0(H)\in\{-1,0\}$; $-1$ is excluded by (b), so $H$'s line set is in the list of \Cref{rot:h1s:CL} (compounds expanded; for a derived structure the degree is that of its line set). Case by case:
\begin{itemize}[nosep,leftmargin=*]
\item $H=(U,\partial f)$ with the far vertex $w$ of $f$ bare, or $H=$ L24, or $H=\{\partial f,{\rm ct}_\sigma\}$ (a split form), or $H=(U,\partial f)$ with the derivative transferred onto $f$ earlier: a bare vertex $w\in H$ has its derivative on a line of $H$, its other scalar field outside $H$ and its photon field outside $H$ (external, or a photon line later than $H$); (R2a) raises $D$ to $1$. The lone $\partial\partial C$ of a split self-energy: (R2a) at a vertex of the split unit (its photon is the unit's $D$ line, outside the one-line component), or, equivalently, (R2b) with $c=\sum_y\partial\partial C(x,y)=0$ exactly (a lattice total derivative), the recipient being the leg at a terminal, which precedes the second unit line in the split regime; then $\{\partial\partial C^{+1},\nabla f\}$ has $D=2$ and $\{\partial\partial C^{+1},\nabla f,D\}$ has $D=1$.
\item $H=T''$ (the triangle's two-point local kernel): the terminals' photon field is the triangle's own $D$ line, inside $T''$, so (R2a) is not available; (R2b) applies with $c=0$ exactly (inversion oddness); $D\mapsto1$.
\item $H=$ W, scalar (outer current vertices bare, derivatives internal): an outer vertex's photon is a line inside W; (R2b) with $c=0$ exactly against the co-switched order-four counterterm share; $D\mapsto1$. $H=$ W, photon, or the double seagull, the sunset, the $e^2\lambda$ members: (R2b) with the common number $c$; $D\mapsto1$. When two triangles inside a W are both bundled (the crossed and vertex-correction graphs: each triangle's outer vertex carries an external leg or a later line, so each qualifies under (R1)), their line sets overlap and the two compounds merge into one element whose line set is all of W: a two-point degree-zero component when it stands alone, handled by (R2b). The collapsed local terms inside W, $\{\partial C,D\}$ and $\{C,\partial C\}$ between the collapsed point and the other terminal, are degree-zero two-point kernels of the SC and VP type: (R2b) with the common number; $D\mapsto1$.
\item Chains $U_1$--$l$--$U_2$, ct--$l$--$U$, ct--$l$--ct with $D=0$: each $U_i$ is bundled on the chain line $l$ or on its outer leg (if neither vertex of $U_i$ qualifies, its outer leg is earlier and attached, and then $H$ contains it with $D\ge1$); if $U_1$ bundles on $l$ and $U_2$ on its outer leg $g$, $\{(U_1,l),(U_2,g)\}$ has $D=(2+d_l)+(2+d_g)\ge1$; if both bundle on $l$, $H$ is the double compound element, which with two units is never alone (its outer legs are earlier) and in chains of three or more units can stand alone (the middle unit forced onto a leg already claimed by a neighbour while both outer legs of the double compound are later): then it is a two-point degree-zero component with no bare vertex and (R2b) applies, with $c=0$ exactly; $D\mapsto1$. A counterterm share stands alone only when its unit is split, so $\{f,g,{\rm ct}\}\ge5+d_f+d_g\ge1$, and $\{\partial f,{\rm ct}\}$ alone never forms since $g$ is earlier too.
\item $H$ one-point (both slots at one vertex: the seagull with two external photons and a hanging self-energy; the all-$C$ components of vacuum monomials cut at their $\Gamma$-lines): a common number, (R2c).
\item The weight classes of the local-term monomials, a weight-$(-1)$ vertex joined by an undifferentiated photon line to a vacuum-polarisation compound whose leg is that line, or to another such vertex ($-1+(-4-1)+3\cdot2=0$, $-1-1-1+3=0$): two-point structures whose outer vertices carry external or later photon fields; (R2b) with the common number; $D\mapsto1$.
\end{itemize}
Nothing else occurs: the list of \Cref{rot:h1s:CL} is exhausted (T and SC are bundled when a vertex qualifies, and otherwise never alone: $\{f_{\rm earlier},T\}$ has $D=d_f+3\ge1$).

(d) \emph{Components with $D=\alpha$.} A compound $(U,E)$ on an external leg alone; any component containing it and a line has $D\ge\alpha+1$.

(e) \emph{Termination and the count.} Each application of (R2) raises one freshly formed component and rewrites one vertex or differences one factor; a vertex is rewritten at most once (after (R2a) at $w$ every later component containing the recipient contains $H_0$ with $D\ge1$); (R2b) is applied once per formed two-point component; the local-term monomials have fewer vertices. The per-flow zero-degree components are the weight-$0$ local vertices and the explicit-$\eta'$ components, at most one per rule application, so $n_0\le\bar n$.
\end{proof}

\begin{remark}[the enumeration]\label{rot:h1s:census}
The case analysis of \Cref{rot:h1s:Dp}(c) is a finite statement for each graph, and an exhaustive enumeration, in exact integer arithmetic, confirms it for every graph with at most two units among the two- and three-point structures through order four (the isolated self-energy and vacuum polarisation, the tree (2.4), the triangle and the seagull--current loop with their leg currents, a unit between two currents, the rainbow, crossed, insertion, vertex-correction and seagull graphs, the chains $U$--$l$--$U$ and $U$--$\partial l$--$w$--$l'$--$U$ with bare and unit vertices, the collapsed local-term graphs and the weight classes): for every derivative placement allowed by \Cref{rot:h1s:cv} and every ordering of the scale-summed lines, walking the initial segments and applying the rule as stated (both monomials of (R2a) followed separately), no component of negative degree forms, every degree-zero component falls under one of the clauses, the rule terminates, and no line carries more than two lattice derivatives; the same holds for the chains of three and four units, the rings and the counterterm monomials, the double compound standing alone in a fraction of the orderings of the three-unit chain and being handled by (R2b). The proof for every family of order $\le\bar n$ is \Cref{rot:h1s:CL} together with (b)--(e), which use only the block arithmetic and the nesting of the initial segments. The statement for chains of any number of units is degree arithmetic: a third unit cannot join a set of bare degree $\le0$ (\Cref{rot:h1s:CL}, $\rho\le2$), and each further unit enters as a compound at a later leg or inside a component of degree $\ge1$.
\end{remark}
